% ============================================================
\section{Iteración 8 --- DRV-08: punto de variación del oponente}\label{sec:it08}
% ============================================================

% ------------------------------------------------------------
\subsection{Driver y alcance}

\begin{tabla}{|L{3.0cm}|Y|}
\fichacab
\parte{Driver}{DRV-08 --- Plazo de 14 semanas con un equipo de dos personas: rebanadas verticales y puntos de variación acotados.}
\parte{Enunciado}{La arquitectura la construye un equipo de dos personas en 14 semanas, por rebanadas verticales que cada semana den un incremento funcional y demostrable, con una regresión sin interfaz, y con el núcleo jugable y el modo en línea antes que lo opcional. Lo que llegará con el proyecto avanzado se aísla en puntos de variación concretos, y solo en ellos.}
\parte{Por qué le toca este turno}{Suma 8 de 20 y es el último con iteración propia porque no cambia qué componentes existen. Lo que queda por decidir son los puntos de variación, y el primero es el más caro y el más probable: el oponente.}
\parte{Qué decide esta iteración}{Cómo se aísla al oponente controlado por IA para que se pueda recortar sin tocar el núcleo, y qué forma tiene un punto de variación en este proyecto.}
\parte{Decisiones ya tomadas}{DEC-01 a DEC-03 y lo decidido en las siete iteraciones anteriores. Pesan CD-01 y CD-02, que ponen al servidor a decidir cada jugada, y CD-22, que da la regresión sin interfaz con la que se comprueba que un recorte no rompió nada.}
\parte{Qué no decide}{Los otros dos puntos de variación: el extremo de red y el modo de despliegue (QAS-14, Q-02) y la marca con los recursos de terceros (QAS-15, Q-07 y CON-18). Tienen la misma forma y quedan para iteraciones posteriores. Tampoco decide Q-01: si habrá IA o no lo deciden STK-04 y STK-06.}
\end{tabla}

% ------------------------------------------------------------
\subsection{Paso 1. Revisión de entradas}

\subsubsection*{ASR que agrupa el driver}

Del conjunto de la sección~\ref{sec:asr}, estos son los ASR que DRV-08 agrupa y el papel que cumple cada uno en esta iteración.

\begin{tabla}{|L{1.4cm}|Y|L{1.7cm}|L{3.6cm}|}
\hline
\filacab \cab{ASR} & \cab{Qué debe cumplir la arquitectura} & \cab{Prioridad} & \cab{Papel en esta iteración} \\ \hline
\endhead
ASR-15 & La solución se construye en 14 semanas, con un incremento auditable cada semana, por un equipo de dos personas, y lo que no es esencial se puede recortar sin tocar lo esencial. & Alta & Objetivo. Su medida es QAS-13. \\ \hline
ASR-10 & Un jugador se enfrenta a un rival controlado por el servidor, con niveles y sugerencias, dentro del mismo presupuesto por jugada. & Media & Condición: la IA existe como un participante que se puede retirar. \\ \hline
ASR-07 & Reservado: partidas en red local con el modo en línea ya entregado. & Reservado & Fuera de esta iteración: solo si Q-02 se resuelve a favor. \\ \hline
\end{tabla}

\subsubsection*{Objetivos de diseño}

\begin{tabla}{|L{1.5cm}|L{5.0cm}|Y|}
\hline
\filacab \cab{ID} & \cab{Objetivo} & \cab{Origen y qué exige aquí} \\ \hline
\endhead
OBJ-24 & El proyecto entrega en 14 semanas, con lo obligatorio terminado. &
CON-13 y CRN-07, con STK-03 y STK-04. Obliga a que lo opcional se pueda soltar. \\ \hline
OBJ-25 & Una decisión que llega tarde no obliga a rehacer lo construido. &
CRN-14, con STK-04 y STK-07. Cuatro preguntas abiertas caerán con el proyecto avanzado. \\ \hline
OBJ-04 & Cada semana hay un incremento funcional y demostrable. &
CON-14, con STK-16. \\ \hline
OBJ-26 & Abstraer solo lo que tiene alta probabilidad de cambiar. &
TEN-04. Cada punto de extensión cuesta tiempo del mismo plazo que protege. \\ \hline
\end{tabla}

\subsubsection*{Requisitos funcionales primarios}

\begin{tabla}{|L{1.3cm}|L{4.2cm}|Y|}
\hline
\filacab \cab{CU} & \cab{Caso de uso} & \cab{Qué exige del punto de variación} \\ \hline
\endhead
CU-27 & Jugar contra la IA &
El rival lo controla el servidor, con cuatro niveles, tablero de 9$\times$9 o 7$\times$7, reloj opcional, deshacer y sugerencias. Es todo lo que habría que recortar, y toca menú, partida y motor. \\ \hline
CU-17, CU-18, CU-19 & Buscar partida clasificatoria; crear sala privada; unirse con código &
Son las otras formas de llegar a una partida. Si la IA desaparece, estas tienen que seguir funcionando sin cambios. \\ \hline
CU-20, CU-21 & Mover el peón; colocar un muro &
El turno se resuelve igual sea quien sea el rival: es lo que permite tratar a la IA como un participante más. \\ \hline
CU-25 & Finalizar la partida y aplicar la progresión &
Distingue las partidas clasificatorias de las que no lo son (RN-03), y las de IA no lo son: la frontera ya existe en el caso de uso. \\ \hline
\end{tabla}

\subsubsection*{Escenarios de atributos de calidad}

\begin{tabla}{|L{1.7cm}|L{3.3cm}|Y|}
\hline
\filacab \cab{Escenario} & \cab{Atributo} & \cab{Medida} \\ \hline
\endhead
QAS-13 & Mantenibilidad: modularidad &
El recorte de la IA en la semana 10 cuesta menos de 1 día de trabajo, se modifican 0 archivos del motor de reglas, del servidor de partidas y de la capa de red, el 100~\% de las pruebas automáticas del núcleo y del modo en línea pasa, el incremento se entrega en su fecha y ninguna opción de la interfaz lleva a una función que ya no existe. \\ \hline
QAS-17 & Mantenibilidad: capacidad de prueba &
Condición. La regresión sin interfaz es lo que permite afirmar que el recorte no rompió nada. \\ \hline
QAS-01 & Desempeño &
Condición. Si la IA existe, su tiempo de cálculo ocupa el mismo presupuesto por jugada que un jugador humano. \\ \hline
\end{tabla}

\subsubsection*{Restricciones}

\begin{tabla}{|L{1.9cm}|L{4.3cm}|Y|}
\hline
\filacab \cab{ID} & \cab{Restricción} & \cab{Cómo limita esta iteración} \\ \hline
\endhead
CON-13 & Plazo máximo de 14 semanas & Es el origen del driver: obliga a que lo adicional sea recortable sin rehacer el núcleo. \\ \hline
CON-14 & Incrementos auditables cada semana & El recorte tiene que caber dentro de un incremento y dejarlo entregable. \\ \hline
REQ-07 & Comprensible para niños de 8 años & La pantalla donde se elige rival no puede quedar con opciones muertas tras el recorte. \\ \hline
CON-01, CON-16 & C\# sobre .NET & La abstracción del participante es una interfaz del mismo ensamblado, de costo bajo. \\ \hline
\end{tabla}

\subsubsection*{Decisiones de proyecto ya tomadas}

\begin{tabla}{|L{1.5cm}|L{4.3cm}|Y|}
\hline
\filacab \cab{ID} & \cab{Decisión} & \cab{Qué impone a esta iteración} \\ \hline
\endhead
DEC-02 & Toda la lógica del juego en el servidor & La IA se ejecuta del lado del servidor, como ya describe CU-27: no puede vivir en el cliente. \\ \hline
CD-04, CD-05 & Máquina de estados y objeto activo por partida & La partida ya trata a sus participantes por igual, lo que abarata tratar a la IA como uno más. \\ \hline
CD-21 & El motor enumera jugadas legales y explica el rechazo & La IA y las sugerencias de CU-27 consumen la misma interfaz del motor. \\ \hline
CD-22 & Arnés de ejecución por lotes sin interfaz & Es la prueba con la que se comprueba que el recorte no rompió el núcleo. \\ \hline
\end{tabla}

\subsubsection*{Concerns}

\begin{tabla}{|L{1.5cm}|Y|}
\hline
\filacab \cab{ID} & \cab{Concern y qué aporta a la iteración} \\ \hline
\endhead
CRN-07 & \emph{Necesito saber si todo lo que se ha comprometido cabe de verdad en las 14 semanas del curso, y qué se recorta primero si no cabe.} \\ \hline
CRN-14 & El director no ha decidido si habrá IA ni conexión LAN, y el equipo no puede quedarse esperando esa decisión. \\ \hline
CRN-08 & Cada semana hay que enseñar algo que funcione y se pueda probar, no una explicación de lo que se construye por dentro. \\ \hline
CRN-03 & La IA es también lo que sostiene el interés de quien ya sabe jugar, así que recortarla tiene un costo que conviene conocer. \\ \hline
\end{tabla}

\subsubsection*{Preguntas abiertas que condicionan la iteración}

\begin{tabla}{|L{1.3cm}|Y|}
\hline
\filacab \cab{ID} & \cab{Cómo condiciona} \\ \hline
\endhead
Q-01 & Si habrá un oponente controlado por IA. Esta iteración no la responde: hace que cualquiera de las dos respuestas sea barata. \\ \hline
Q-02, Q-07, Q-08 & Los otros puntos de variación y las variantes de reglas. Comparten forma con este, y las variantes ya entran como parámetros por CD-19. \\ \hline
\end{tabla}

% ------------------------------------------------------------
\subsection{Paso 2. Objetivo de la iteración y selección de entradas}

\subsubsection*{Priorización de las entradas}

\begin{tabla}{|L{3.2cm}|L{1.3cm}|Y|L{2.2cm}|}
\hline
\filacab \cab{Entrada} & \cab{Peso} & \cab{Por qué pesa lo que pesa} & \cab{Decisión} \\ \hline
\endhead
QAS-13 & Alto & Es la medida del objetivo y fija el tope de costo del recorte: un día y 0 archivos del núcleo. & Se atiende \\ \hline
CU-27 & Alto & Es todo lo que habría que recortar, y describe cuánto alcanza: menú, partida, motor y progresión. & Se atiende \\ \hline
CRN-07, CRN-14 & Alto & Son el origen del driver: el alcance puede no caber y la decisión llega tarde. & Se atiende \\ \hline
TEN-04 & Alto & Limita la solución: la abstracción tiene que ser barata, o el remedio consume el plazo que protege. & Se atiende \\ \hline
CON-13, CON-14 & Alto & Fijan el plazo y la forma de entrega en la que el recorte tiene que caber. & Se atiende \\ \hline
CU-17, CU-18, CU-19, CU-25 & Medio & Definen lo que tiene que seguir funcionando igual después del recorte. & Se atiende como condición \\ \hline
QAS-17, CD-22 & Medio & Son la evidencia de que nada se rompió; ya existen y aquí se usan. & Se atiende como supuesto \\ \hline
QAS-01 & Medio, como límite & Si la IA existe, su cálculo entra en el presupuesto por jugada ya repartido. & Se atiende como restricción \\ \hline
Q-01 & Bajo & No se resuelve: la iteración vuelve baratas las dos respuestas. & Se atiende como condición \\ \hline
QAS-14, QAS-15 & Bajo aquí & Son los otros dos puntos de variación, con la misma forma y otro elemento. & Se difieren a iteraciones posteriores \\ \hline
\end{tabla}

\subsubsection*{Objetivo de la iteración}

\begin{tabla}{|L{3.0cm}|Y|}
\fichacab
\parte{Objetivo}{Que el oponente controlado por IA se pueda retirar del alcance sin tocar el núcleo jugable ni el modo en línea.}
\parte{ASR que satisface}{ASR-15, con ASR-10 como condición.}
\parte{Driver que cubre}{DRV-08, en su primer punto de variación.}
\parte{Medida}{La de QAS-13: el recorte cuesta menos de 1 día de trabajo, se modifican 0 archivos del motor de reglas, del servidor de partidas y de la capa de red, el 100~\% de las pruebas automáticas del núcleo y del modo en línea pasa, el incremento de esa semana se entrega en su fecha y ninguna opción de la interfaz lleva a una pantalla rota.}
\parte{Límite que respeta}{TEN-04 y QAS-01. La abstracción tiene que costar poco de construir, y si la IA se queda, su tiempo de cálculo cabe en el presupuesto por jugada de la iteración~4.}
\parte{Incremento que produce}{Una partida contra la IA jugada de principio a fin, y en una rama aparte el ensayo del recorte, con el registro de las horas empleadas, la lista de archivos modificados y la regresión ejecutada. La demostración es doble: el juego funciona con IA y funciona igual sin ella.}
\parte{Criterio de éxito}{La iteración termina cuando el ensayo del recorte cumple las cinco cifras de la medida.}
\end{tabla}

% ------------------------------------------------------------
\subsection{Paso 3. Elemento a descomponer}

% ------------------------------------------------------------
\Needspace*{0.4\textheight}
\subsubsection*{EL-08 --- El punto de variación del oponente}

\begin{tabla}{|L{2.6cm}|Y|}
\fichacab
\parte{Elemento}{El punto de variación del oponente: el lugar del sistema donde se decide quién ocupa el asiento del rival, una persona conectada o un jugador controlado por el servidor.}
\parte{Qué abarca}{La forma en que la partida se dirige a cada participante para pedirle su jugada y notificarle el estado, el módulo que juega por la IA con sus cuatro niveles y sus sugerencias (CU-27), la pantalla donde el jugador elige contra quién juega y la regla que decide qué partidas son clasificatorias (RN-03 de CU-25). No abarca el motor, que es el mismo para todos, ni el arbitraje, que no distingue entre rivales.}
\parte{Por qué se elige}{Porque la medida del objetivo describe qué pasa cuando este punto se apaga: cuánto cuesta, qué archivos se tocan y qué sigue funcionando. Se elige frente a descomponer el módulo de IA, que es lo que se recorta pero no lo que hace que el recorte sea barato, y frente al servicio de estado de la partida, que precisamente no debe cambiar. Es el primero de los tres puntos de variación porque es el más caro y el que CRN-07 nombra como primer candidato al recorte.}
\parte{Qué falta decidir sobre él}{Cómo se dirige la partida a un participante que no está al otro lado de una conexión. Dónde se decide que una modalidad existe o no, y cómo desaparece de la interfaz sin dejar opciones muertas. Y hasta dónde llega la abstracción, porque TEN-04 impide abstraer más de lo que tiene alta probabilidad de cambiar.}
\parte{Evidencia en los casos de uso}{CU-27 dice que el rival lo controla el servidor, lo que confirma que la IA vive donde vive la autoridad, y describe ajustes, niveles, sugerencias y deshacer, que es alcance que solo existe en esa modalidad. CU-25 ya separa las partidas clasificatorias de las demás en su RN-03, así que la frontera de la progresión está trazada. CU-20 y CU-21 no distinguen quién es el rival: resuelven el turno igual, que es lo que hace posible tratar a la IA como un participante más.}
\parte{Trazabilidad}{DRV-08. DEC-02. CD-01, CD-04, CD-05, CD-21, CD-22. QAS-13, con QAS-17 y QAS-01 como condiciones. CU-27, CU-17, CU-18, CU-19, CU-20, CU-21, CU-25 (RN-03). CON-13, CON-14, REQ-07. CRN-07, CRN-14, CRN-08, CRN-03. TEN-04. Q-01. STK-04, STK-03, STK-07, STK-06.}
\end{tabla}

% ------------------------------------------------------------
\subsection{Paso 4. Conceptos de diseño}

% ------------------------------------------------------------
\Needspace*{0.3\textheight}
\subsubsection*{CD-45 --- El participante como abstracción única}

\begin{tabla}{|L{2.6cm}|Y|}
\fichacab
\parte{Estilo}{Patrón de diseño: interfaz común para variantes.}
\parte{Descripción}{La partida se dirige a sus participantes a través de una sola abstracción, que recibe el estado y entrega una jugada, y es indiferente a que detrás haya una conexión o un módulo que calcula. Atiende QAS-13 en su parte más exigente, 0 archivos modificados en el arbitraje: si la partida no sabe qué tipo de rival tiene, retirar uno no la toca. Se elige frente a que el arbitraje pregunte de qué tipo es el rival en cada paso, que repartiría esa condición por todo el servicio y haría del recorte una cacería de condicionales.}
\parte{Efectos}{El módulo de IA se convierte en un consumidor más del motor (CD-21) y del estado, sin acceso privilegiado. Deja preparado el caso de CU-27 sin decidir Q-01. La abstracción tiene que admitir que un participante tarde en responder, lo que ya ocurre con una persona y encaja con el presupuesto de CD-23.}
\parte{Trazabilidad}{DRV-08, DRV-01. DEC-02. CD-04, CD-05, CD-21. QAS-13, QAS-01. CU-27, CU-20, CU-21. CON-13, CON-01. CRN-14, CRN-07. TEN-04. Q-01. STK-05, STK-07, STK-04.}
\end{tabla}

% ------------------------------------------------------------
\Needspace*{0.3\textheight}
\subsubsection*{CD-46 --- Las modalidades se declaran en configuración}

\begin{tabla}{|L{2.6cm}|Y|}
\fichacab
\parte{Estilo}{Patrón de configuración, con táctica de modularidad.}
\parte{Descripción}{Las modalidades disponibles, en línea, sala privada y contra la IA, se declaran en configuración, y la interfaz muestra solo las declaradas. Atiende la última parte de QAS-13, que ninguna opción lleve a una función que ya no existe, y hace que el recorte empiece por un cambio de configuración y no por borrar código. Se elige frente a escribir las opciones en la pantalla, que obliga a modificar la interfaz en cada recorte y a arriesgarse a dejar botones muertos.}
\parte{Efectos}{El recorte se ensaya en minutos antes de decidirlo, lo que da a STK-04 información real para elegir. La misma declaración sirve para apagar una modalidad en producción si algo falla. Obliga a que la interfaz se construya a partir de la lista y no con opciones fijas.}
\parte{Trazabilidad}{DRV-08. CD-19, CD-45. QAS-13, QAS-02. CU-27, CU-17, CU-18. CON-13, CON-14, REQ-07. CRN-07, CRN-14. TEN-04. Q-01, Q-02. STK-04, STK-09, STK-01.}
\end{tabla}

% ------------------------------------------------------------
\Needspace*{0.3\textheight}
\subsubsection*{CD-47 --- Lo opcional depende del núcleo, nunca al revés}

\begin{tabla}{|L{2.6cm}|Y|}
\fichacab
\parte{Estilo}{Estilo: dirección de las dependencias.}
\parte{Descripción}{El módulo de IA depende del motor, del estado y de la abstracción de participante, y ninguno de ellos depende de él. Atiende QAS-13 en la parte de los 0 archivos modificados en el núcleo: lo que nadie usa se puede quitar sin dejar huecos. Se elige frente a dejar que el núcleo llame a la IA para pedirle sugerencias o consejos, que convertiría el recorte en una modificación del núcleo, que es justo lo que la medida prohíbe.}
\parte{Efectos}{Las sugerencias de CU-27 se resuelven pidiendo al motor las jugadas legales y evaluándolas dentro del módulo de IA, no al revés. La regla se aplica igual a los otros dos puntos de variación pendientes, así que fija la forma de todos. Obliga a revisar la dirección de cada dependencia nueva en la revisión semanal.}
\parte{Trazabilidad}{DRV-08, DRV-05. CD-18, CD-21, CD-45. QAS-13, QAS-17. CU-27, CU-25 (RN-03). CON-13, CON-14. CRN-07, CRN-14, CRN-20, CRN-24. TEN-04. Q-01, Q-02, Q-07. STK-05, STK-07, STK-04.}
\end{tabla}

% ------------------------------------------------------------
\Needspace*{0.3\textheight}
\subsubsection*{CD-48 --- Rebanadas verticales con la regresión como cierre}

\begin{tabla}{|L{2.6cm}|Y|}
\fichacab
\parte{Estilo}{Táctica de proceso con efecto estructural: orden de construcción.}
\parte{Descripción}{Cada semana se construye una rebanada que atraviesa interfaz, protocolo, arbitraje y datos hasta dar algo jugable, y el incremento se cierra ejecutando la regresión sin interfaz de CD-22 junto con la medida del escenario que tocó esa semana. Atiende QAS-13 en la parte de entregar el incremento en su fecha y en la de demostrar que el núcleo y el modo en línea siguen funcionando. Se elige frente a construir capa por capa, que deja lo demostrable para el final y hace que un recorte tardío no tenga con qué compararse.}
\parte{Efectos}{Las medidas de los escenarios se convierten en la batería de pruebas del proyecto, que es la evidencia que STK-16 audita cada semana.}
\parte{Trazabilidad}{DRV-08, DRV-05. CD-22, CD-27, CD-44. QAS-13, QAS-17. CU-20, CU-21, CU-27. CON-13, CON-14. CRN-07, CRN-08, CRN-17. TEN-04. STK-03, STK-16, STK-04, STK-10.}
\end{tabla}
